\documentclass[10pt,a4paper]{amsart}
\usepackage[margin=19mm]{geometry}
\usepackage{amsmath,amssymb,mathtools}
\usepackage[scr=rsfs,cal=euler]{mathalfa}
\usepackage{xfrac}
\usepackage[unicode,hidelinks,bookmarks=false]{hyperref}
\newcommand{\ie}{i.e.\ }
\newcommand{\cf}{cf.\ }
\newcommand{\resp}{respectively\ }
\newcommand{\Subsection}{\subsection}
\providecommand{\nobreakdash}{\nobreak\hbox{-}}
\setlength{\emergencystretch}{2em}
\multlinegap0pt
\usepackage{xcolor}
\usepackage{amssymb}
\usepackage{enumerate}
\usepackage{float}
\usepackage{tikz}
\usepackage{mathtools}
\usepackage{bm}
\usepackage{tcolorbox}
\newtheorem{lemma}{Lemma}
\newcommand{\C}{\ensuremath{\mathbb{C}}} %complex numbers
\newcommand{\DPS}{\text{\rm DPS}}
\newcommand{\Herm}{\mathrm{Herm}}
\newcommand{\ML}{\textcolor{blue}}
\newcommand{\Perm}{\mathfrak G}
\newcommand{\Tr}{\text{\rm Tr}}
\newcommand{\ot}{\otimes}
\newcommand{\uh}{\underline{h}}
\newcommand{\ui}{\underline{i}}
\newcommand{\uj}{\underline{j}}
\newcommand{\uk}{\underline{k}}
\title{Semidefinite hierarchies for diagonal unitary invariant bipartite quantum states}
\author{Jonas Britz, Monique Laurent}
\date{Source: arXiv:2512.06551v2 (CC BY 4.0)}
\begin{document}
\maketitle

\noindent\emph{Source and licence.} Semidefinite hierarchies for diagonal unitary invariant bipartite quantum states. Version arXiv:2512.06551v2 (CC BY 4.0), distributed by arXiv under the Creative Commons Attribution 4.0 International licence, http://creativecommons.org/licenses/by/4.0/, as recorded in the arXiv licence metadata for this exact version and shown on the abstract page https://arxiv.org/abs/2512.06551v2. Full licence notice: \texttt{licenses/arxiv-2512.06551-CC-BY-4.0.txt}.\par\smallskip

\noindent\textit{Editorial scope.} Counted unit: the article's lemma LThetaAppliedDPSLCertificate (source0.tex lines 1869--1875). The article's complete original proof of it is delivered with it (source0.tex lines 1877--1916, one \texttt{proof} environment of 40 lines). No mathematics is written, repaired or re-derived: the statement and the proof are the article's own lines, in the article's own notation, and no retained line is substituted or re-flowed.  Retained uncounted as the source-local context the proof needs: lines 664--669; lines 672--674; lines 679--681; lines 1855--1857.  Nothing was omitted from any retained span, so the package carries no omission record.  No retained line contains a citation, so the wrapper carries no bibliography and no bibliography entry.  No retained line contains a figure environment, an image-inclusion command or drawing code, so no image is delivered and the package declares no asset dependency.  Editorial formatting only, in addition: an AMS \texttt{amsart} preamble that restates the article's own declarations and macros used by the retained lines, namely the environments \texttt{lemma} on separate counters, together with the standard packages those lines need.  The retained environments print in this document as Lemma 1 in order of appearance, whereas the article prints the same items with the numbering of its own full document.  The complete native source of the article as distributed in the arXiv e-print for this exact version is bundled unchanged as \texttt{sources/190504/source0.tex}.
\begin{align}\label{eq:DPS1}
\begin{split}
& (I_n\ot \Pi_t) \ \rho_{AB[t]} \ (I_n\ot \Pi_t) = \rho_{AB[t]}; \text{  equivalently}, \  \rho_{AB[t]} \ (I_n\ot \Pi_t) = \rho_{AB[t]}, \\
&\text{i.e., } (\rho_{AB[t]})_{i_0\ui,j_0\uj}= (\rho_{AB[t]})_{i_0 \ui,j_0\uj^\sigma} \text{ for all } \sigma \in \Perm_t,\ i_0,j_0\in [n],\ \ui,\uj\in [n]^t;\\
& \text{in other words, } \rho_{AB[t]}\in \Herm(\C^n\ot S^t(\C^n)).
\end{split}\end{align}

\begin{align}\label{eq:DPS2}
\rho_{AB[t]}^{T_{B[s]}}\succeq 0  \ \text{ for all } s=0,1,\ldots,t.
\end{align}

\begin{align}\label{eq:DPS3}
\Tr_{B[2:t]}(\rho_{AB[t]}) = \rho_{AB}.
\end{align}

\begin{align}\label{eq:theta-sigma}
\theta(\ui^\sigma) =\theta(i_{\sigma(1)})\ldots\theta(i_{\sigma(t)})= (\theta(i_1)\ldots\theta(i_t))^\sigma=\theta(\ui)^\sigma \text{ for any } \theta\in\Perm_n, \sigma\in \Perm_t.
\end{align}

\begin{lemma}\label{LThetaAppliedDPSLCertificate}
 Assume $\rho_{AB}\in \DPS^{(t)}_n$ with $\rho_{AB[t]}$ as $\DPS^{(t)}$-certificate.
 %\footnote{\ML{It does not seem that the condition $\rho_{AB}=\theta(\rho_{AB})$ is needed here.}}
For any $\theta\in \Perm_n$,  $\theta(\rho_{AB[t]}) $ is a $\DPS^{(t)}$-certificate for $\theta(\rho_{AB})$.
%, which implies $\theta(\rho_{AB})\in\DPS^{(t)}_n$.
 %    Let $\theta\in Perm(n)$, $\rho_{AB}\in\DPS^{(t)}$ with certificate $\rho_{AB[t]}$ and $\rho_{AB}^\theta=\rho_{AB}$. Then $\rho_{AB[t]}^\theta$ is also a certificate. 
\end{lemma}

\begin{proof}
  By assumption,   $\rho_{AB[t]}$ satisfies (\ref{eq:DPS1}),  (\ref{eq:DPS2}), (\ref{eq:DPS3}), we show that the same holds for $\theta(\rho_{AB[t]})$.
First, we show (\ref{eq:DPS1}). For this,    let $\sigma\in \Perm_t$ and $i_0\ui,j_0\uj\in [n]^{t+1}$. Then, we have
    \begin{align*}
        \theta(\rho_{AB[t]})_{i_0\ui^\sigma,j_0\uj^\sigma}&=(\rho_{AB[t]})_{\theta(i_0)\theta(\ui^\sigma),\theta(j_0)\theta(\uj^\sigma)}
        =(\rho_{AB[t]})_{\theta(i_0)(\theta(\ui))^\sigma,\theta(j_0)\theta(\uj)^\sigma}\\
        &=(\rho_{AB[t]})_{\theta(i_0)\theta(\ui),\theta(j_0)\theta(\uj)}
        =\theta(\rho_{AB[t]})_{i_0\ui,j_0\uj},
    \end{align*}
using (\ref{eq:theta-sigma}) at the first line and the fact that $\rho_{AB[t]}$ satisfies \eqref{eq:DPS1} at the second line. Hence, $\theta(\rho_{AB[t]})$ also satisfies \eqref{eq:DPS1}.

    Condition \eqref{eq:DPS2} follows from the fact that the operations of applying $\theta$ and taking the partial transpose commute:
  %  \begin{align*}
       $ \left(\theta(\rho_{AB[t]})\right)^{T_{B[s]}}=\theta\big(\rho_{AB[t]}^{T_{B[s]}}\big),$
%    \end{align*}
 which can be checked by computing
    \begin{align*}
    &    \big( \left(\theta(\rho_{AB[t]})\right)^{T_{B[s]}}\big)_{i_0\ui,j_0\uj} =
        \left(\theta(\rho_{AB[t]})\right)_{i_0j_1\dots j_si_{s+1}\dots s_t,j_0i_1\dots i_sj_{s+1}\dots j_t}\\
        &=(\rho_{AB[t]})_{\theta(i_0j_1\dots j_si_{s+1}\dots s_t),\theta(j_0i_1\dots i_sj_{s+1}\dots j_t)}
        =\big(\rho_{AB[t]}^{T_{B[s]}}\big)_{\theta(i_0\ui),\theta(j_0\uj)}\\
       & =\big(\theta\big(\rho_{AB[t]}^{T_{B[s]}}\big)\big)_{i_0\ui,j_0\uj}.
    \end{align*}
 Hence, $\rho_{AB[t]}^{T_{B[s]}}\succeq 0$ implies $\theta(\rho_{AB[t]}^{T_{B[s]}})\succeq 0$ and thus
  $(\theta(\rho_{AB[t]}))^{T_{B[s]}}\succeq 0$, as desired. 
      
       Finally, condition \eqref{eq:DPS3} follows from the fact that the operations of applying $\theta$ and taking the partial trace also commute, which 
       can be checked   by computing
    \begin{align*}
   &     \big(\Tr_{B[2:t]}\big(\theta(\rho_{AB[t]})\big)\big)_{i_0i_1,j_0j_1}
%        &=\sum_{k_2,\dots,k_t=1}^{n}\big(\theta(\rho_{AB[t]})\big)_{i_0i_1k_2\dots k_t,j_0j_1k_2\dots k_t}\\
  =\sum_{\uk\in [n]^{t-1}} \big(\theta(\rho_{AB[t]})\big)_{i_0i_1\uk, j_0j_1\uk}\\
 %       &=\sum_{k_2,\dots,k_t=1}^{n}\big(\rho_{AB[t]}\big)_{\theta(i_0i_1k_2\dots k_t),\theta(j_0j_1k_2\dots k_t)}\\
  &= \sum_{\uk\in [n]^{t-1}} \big(\rho_{AB[t]}\big)_{\theta(i_0i_1)\theta(\uk),\theta(j_0j_1)\theta(\uk)}
  %       &=\sum_{h_2,\dots,h_t=1}^{n}\big(\rho_{AB[t]}\big)_{\theta(i_0i_1)h_2\dots h_t,\theta(j_0j_1)h_2\dots h_t}\\
  =\sum_{\uh\in [n]^{t-1}} \big(\rho_{AB[t]}\big)_{\theta(i_0i_1)\uh, \theta(j_0j_1)\uh} \\
    &      =\Tr_{r-2}(\rho_{AB[t]})_{\theta(i_0i_1),\theta(j_0j_1)}=(\rho_{AB})_{\theta(i_0i_1),\theta(j_0j_1)} 
           =\theta(\rho_{AB})_{i_0i_1,j_0j_1}, %=(\rho_{AB})_{i_0i_1,j_0j_1}
    \end{align*}
using the fact that $\rho_{AB[t]}$ satisfies (\ref{eq:DPS3}) at the last but one equality.\qed \end{proof}

\end{document}
